% GENERATED FILE. Edit canonical lessons, then run npm run generate:books.
% canonical-source-hash: fnv1a64:4581b8edc5301dfa
% canonical-lessons: MW-C245-terno, MW-C245-juhar, MW-C245-dhukno, MW-C245-achambho, MW-C245-badhai

\chapter{Neighbours Come to Call}
\label{ch:mw-neighbours-come-to-call}
\hlchaptermodality{245}

\begin{chapteropening}
\textbf{By the end of this chapter:} \emph{I can say who came and who greeted whom.}

Complete the last lesson of chapter 245: I can say who came and who greeted whom.
\end{chapteropening}

\section[\texorpdfstring{\mw{तेड़णो} (teṛṇo)}{तेड़णो (teṛṇo)}]{\mw{तेड़णो} (teṛṇo) — to call someone over; to invite}

\label{lesson:MW-C245-terno}



\begin{quote}
Before the new one: say the Marwadi for a fault, then the Marwadi for a threat.
\end{quote}

\subsection*{You'll want to know: \mw{तेड़णो}}
\textbf{\mw{तेड़णो}} — \emph{teṛṇo} — \textquotedblleft{}to call someone over; to invite\textquotedblright{}.

\begin{quote}
\mw{मां} \mw{पावणा} \mw{तेड़ै} \mw{है।}
\end{quote}

\begin{quote}
\emph{mā̃ pāvṇā teṛai hai.}
\end{quote}

\begin{quote}
Mother invites guests over.
\end{quote}

Past (I …): \mw{म्हैं} \mw{पावणा} \mw{तेड़्या}, \emph{mhaĩ pāvṇā teṛyā}.

\begin{cousinweb}
\textbf{\mw{तेड़णो}} (\emph{teṛṇo}) is inherited from a reconstructed \emph{*tēḍati}, \textquotedblleft{}summons\textquotedblright{}; CDIAL lists an Old Marwari noun \emph{teṛū̃}, \textquotedblleft{}a summons\textquotedblright{}, under it (CDIAL 5947a).
\end{cousinweb}

\begin{grammarlens}[title={What you've built}]
One more word for saying what happened.
\end{grammarlens}

\subsection*{Your turn}
\begin{itemize}
\raggedright
  \item \emph{Say it:} \emph{teṛṇo}
  \item \emph{Say it:} \emph{teṛṇo}, once more
  \item \emph{Say it:} \emph{mā̃ pāvṇā teṛai hai.} — Mother invites guests over.
  \item \emph{Recall:} say the Marwadi for a fault, then the Marwadi for a threat, then say \emph{teṛṇo} again
\end{itemize}

\begin{tcolorbox}[breakable,colback=teal!4,colframe=teal!35!black,title={Before you move on}]
What does \mw{तेड़णो} mean? (\textquotedblleft{}To call someone over; to invite\textquotedblright{}.) Say it once more.
\end{tcolorbox}
\section[\texorpdfstring{\mw{जुहार} (juhār)}{जुहार (juhār)}]{\mw{जुहार} (juhār) — a respectful greeting}

\label{lesson:MW-C245-juhar}



\begin{quote}
Before the new one: say the Marwadi for a threat, then the Marwadi for to call someone over.
\end{quote}

\subsection*{You'll want to know: \mw{जुहार}}
\textbf{\mw{जुहार}} — \emph{juhār} — \textquotedblleft{}a respectful greeting\textquotedblright{}.

\begin{quote}
\mw{म्हैं} \mw{जुहार} \mw{कियो।}
\end{quote}

\begin{quote}
\emph{mhaĩ juhār kiyo.}
\end{quote}

\begin{quote}
I greeted them respectfully.
\end{quote}

A masculine noun.

\begin{cousinweb}
\textbf{\mw{जुहार}} (\emph{juhār}) is inherited from Sanskrit \emph{jyotkāra}, \textquotedblleft{}a shout of acclamation\textquotedblright{}; CDIAL lists Old Marwari \emph{juhāra}, a greeting said with the words \textquotedblleft{}live long\textquotedblright{}, under it (CDIAL 5300a).
\end{cousinweb}

\begin{grammarlens}[title={What you've built}]
One more word for saying what happened.
\end{grammarlens}

\subsection*{Your turn}
\begin{itemize}
\raggedright
  \item \emph{Say it:} \emph{juhār}
  \item \emph{Say it:} \emph{juhār}, once more
  \item \emph{Say it:} \emph{mhaĩ juhār kiyo.} — I greeted them respectfully.
  \item \emph{Recall:} say the Marwadi for a threat, then the Marwadi for to call someone over, then say \emph{juhār} again
\end{itemize}

\begin{tcolorbox}[breakable,colback=teal!4,colframe=teal!35!black,title={Before you move on}]
What does \mw{जुहार} mean? (\textquotedblleft{}A respectful greeting\textquotedblright{}.) Say it once more.
\end{tcolorbox}
\section[\texorpdfstring{\mw{ढूकणो} (ḍhūkṇo)}{ढूकणो (ḍhūkṇo)}]{\mw{ढूकणो} (ḍhūkṇo) — to draw near; to arrive at the door}

\label{lesson:MW-C245-dhukno}



\begin{quote}
Before the new one: say the Marwadi for to call someone over, then the Marwadi for a respectful greeting.
\end{quote}

\subsection*{You'll want to know: \mw{ढूकणो}}
\textbf{\mw{ढूकणो}} — \emph{ḍhūkṇo} — \textquotedblleft{}to draw near; to arrive at the door\textquotedblright{}.

\begin{quote}
\mw{पावणो} \mw{घर} \mw{कनै} \mw{ढूक्यो।}
\end{quote}

\begin{quote}
\emph{pāvṇo ghar kanai ḍhūkyo.}
\end{quote}

\begin{quote}
The guest drew near the house.
\end{quote}

Past (he …): \mw{वो} \mw{पावणो} \mw{ढूक्यो}, \emph{vo pāvṇo ḍhūkyo}.

\begin{cousinweb}
\textbf{\mw{ढूकणो}} (\emph{ḍhūkṇo}) is inherited from a reconstructed \emph{*ḍhukyati}, \textquotedblleft{}approaches\textquotedblright{}; CDIAL lists the Old Marwari partner \emph{ḍhukāvaï}, \textquotedblleft{}brings near\textquotedblright{}, and Hindi \emph{ḍhūknā}, \textquotedblleft{}to approach\textquotedblright{}, under it (CDIAL 5592).
\end{cousinweb}

\begin{grammarlens}[title={What you've built}]
One more word for saying what happened.
\end{grammarlens}

\subsection*{Your turn}
\begin{itemize}
\raggedright
  \item \emph{Say it:} \emph{ḍhūkṇo}
  \item \emph{Say it:} \emph{ḍhūkṇo}, once more
  \item \emph{Say it:} \emph{pāvṇo ghar kanai ḍhūkyo.} — The guest drew near the house.
  \item \emph{Recall:} say the Marwadi for to call someone over, then the Marwadi for a respectful greeting, then say \emph{ḍhūkṇo} again
\end{itemize}

\begin{tcolorbox}[breakable,colback=teal!4,colframe=teal!35!black,title={Before you move on}]
What does \mw{ढूकणो} mean? (\textquotedblleft{}To draw near; to arrive at the door\textquotedblright{}.) Say it once more.
\end{tcolorbox}
\section[\texorpdfstring{\mw{अचंभो} (achambho)}{अचंभो (achambho)}]{\mw{अचंभो} (achambho) — a surprise; a wonder}

\label{lesson:MW-C245-achambho}



\begin{quote}
Before the new one: say the Marwadi for a respectful greeting, then the Marwadi for to draw near.
\end{quote}

\subsection*{You'll want to know: \mw{अचंभो}}
\textbf{\mw{अचंभो}} — \emph{achambho} — \textquotedblleft{}a surprise; a wonder\textquotedblright{}.

\begin{quote}
\mw{मामो} \mw{आयो}! \mw{यो} \mw{अचंभो} \mw{है।}
\end{quote}

\begin{quote}
\emph{māmo āyo! yo achambho hai.}
\end{quote}

\begin{quote}
My uncle has come! What a surprise.
\end{quote}

A masculine noun.

\begin{cousinweb}
\textbf{\mw{अचंभो}} (\emph{achambho}) is inherited from a reconstructed \emph{*ācchambha}, \textquotedblleft{}rigidity\textquotedblright{}, and so \textquotedblleft{}astonishment\textquotedblright{}, the stiffness of surprise. CDIAL gives no Marwari form in this entry; it lists Hindi \emph{acambhā}, \textquotedblleft{}wonder\textquotedblright{} (CDIAL 1079).
\end{cousinweb}

\begin{grammarlens}[title={What you've built}]
One more word for saying what happened.
\end{grammarlens}

\subsection*{Your turn}
\begin{itemize}
\raggedright
  \item \emph{Say it:} \emph{achambho}
  \item \emph{Say it:} \emph{achambho}, once more
  \item \emph{Say it:} \emph{māmo āyo! yo achambho hai.} — My uncle has come! What a surprise.
  \item \emph{Recall:} say the Marwadi for a respectful greeting, then the Marwadi for to draw near, then say \emph{achambho} again
\end{itemize}

\begin{tcolorbox}[breakable,colback=teal!4,colframe=teal!35!black,title={Before you move on}]
What does \mw{अचंभो} mean? (\textquotedblleft{}A surprise; a wonder\textquotedblright{}.) Say it once more.
\end{tcolorbox}
\section[\texorpdfstring{\mw{बधाई} (badhāī)}{बधाई (badhāī)}]{\mw{बधाई} (badhāī) — congratulations; good wishes}

\label{lesson:MW-C245-badhai}



\begin{quote}
Before the new one: say the Marwadi for to draw near, then the Marwadi for a surprise.
\end{quote}

\subsection*{You'll want to know: \mw{बधाई}}
\textbf{\mw{बधाई}} — \emph{badhāī} — \textquotedblleft{}congratulations; good wishes\textquotedblright{}.

\begin{quote}
\mw{पाड़ोसी} \mw{बधाई} \mw{देवै} \mw{है।}
\end{quote}

\begin{quote}
\emph{pāṛosī badhāī devai hai.}
\end{quote}

\begin{quote}
The neighbour offers congratulations.
\end{quote}

A feminine noun.

\begin{cousinweb}
\textbf{\mw{बधाई}} (\emph{badhāī}) is inherited from Sanskrit \emph{vardhāpayati}, \textquotedblleft{}gladdens, congratulates\textquotedblright{}; CDIAL lists Old Marwari \emph{badhāī}, \textquotedblleft{}a present\textquotedblright{}, and Hindi \emph{badhāī}, \textquotedblleft{}congratulations\textquotedblright{}, under it (CDIAL 11383).
\end{cousinweb}

\begin{grammarlens}[title={What you've built}]
One more word for saying what happened.
\end{grammarlens}

\subsection*{Your turn}
\begin{itemize}
\raggedright
  \item \emph{Say it:} \emph{badhāī}
  \item \emph{Say it:} \emph{badhāī}, once more
  \item \emph{Say it:} \emph{pāṛosī badhāī devai hai.} — The neighbour offers congratulations.
  \item \emph{Recall:} say the Marwadi for to draw near, then the Marwadi for a surprise, then say \emph{badhāī} again
\end{itemize}

\begin{tcolorbox}[breakable,colback=teal!4,colframe=teal!35!black,title={Before you move on}]
What does \mw{बधाई} mean? (\textquotedblleft{}Congratulations; good wishes\textquotedblright{}.) Say it once more.
\end{tcolorbox}
